\documentclass[11pt]{article}

\usepackage[margin=1in]{geometry}
\usepackage{amsmath,amssymb,amsthm}
\usepackage{graphicx}
\usepackage[hidelinks]{hyperref}

\newtheorem{theorem}{Theorem}
\newtheorem{proposition}{Proposition}

\title{A Reflexive Case of the \texorpdfstring{$L^p$}{Lp} State-Space Problem}
\author{Automatic Math Discovery Smoke Test}
\date{June 9, 2026}

\begin{document}
\maketitle

\section*{Source problem}

The source paper asks for an analogue, for $1<p<\infty$, of its
characterization of weakly compact state spaces in $L^1(\mu,X)$.

\begin{center}
\includegraphics[width=\textwidth]{figures/open_problem_crop.png}
\end{center}

\section*{Result}

This packet gives the full reflexive-domain answer and shows that the
$L^1$ singleton phenomenon cannot extend to $p>1$.

\begin{theorem}
Let $(\Omega,\Sigma,\mu)$ be a finite measure space, let $1<p<\infty$, and let
$q$ be conjugate to $p$. If $X$ is reflexive, then for every
$f\in S_{L^p(\mu,X)}$ the state space
\[
S_f=\left\{g\in B_{L^q(\mu,X^*)}:
  \int_\Omega g(\omega)(f(\omega))\,d\mu(\omega)=1\right\}
\]
is weakly compact in $L^q(\mu,X^*)$.
\end{theorem}

\begin{proof}
Since $X$ is reflexive, $X^*$ is reflexive. Because $1<q<\infty$ and
$\mu$ is finite, the Bochner space $L^q(\mu,X^*)$ is reflexive. Define
\[
\Phi_f(g)=\int_\Omega g(\omega)(f(\omega))\,d\mu(\omega),
        \qquad g\in L^q(\mu,X^*).
\]
Holder's inequality gives $\|\Phi_f\|\leq \|f\|_p=1$, and the usual duality
formula gives $\|\Phi_f\|=1$. Hence
\[
S_f=B_{L^q(\mu,X^*)}\cap \Phi_f^{-1}(\{1\}).
\]
This is a norm-closed convex bounded subset of the reflexive space
$L^q(\mu,X^*)$. Closed convex subsets are weakly closed, so $S_f$ is weakly
closed and bounded in a reflexive space. Therefore $S_f$ is weakly compact.
\end{proof}

\begin{proposition}
For $p>1$, weak compactness of $S_f$ does not imply that $S_f$ is a singleton,
even when $X$ is two-dimensional.
\end{proposition}

\begin{proof}
Take $\Omega=[0,1]$ with Lebesgue measure and
$X=(\mathbb R^2,\|\cdot\|_1)$, so $X^*=(\mathbb R^2,\|\cdot\|_\infty)$.
Let $f(\omega)=e_1=(1,0)$ for all $\omega$. Then $\|f\|_p=1$.

If $h\in L^\infty[0,1]$ satisfies $|h|\leq 1$ almost everywhere, set
\[
g_h(\omega)=(1,h(\omega))\in X^*.
\]
Then
\[
\|g_h\|_{L^q(\mu,X^*)}
=\left(\int_0^1 \max\{1,|h(\omega)|\}^q\,d\omega\right)^{1/q}=1
\]
and
\[
\int_0^1 g_h(\omega)(f(\omega))\,d\omega
=\int_0^1 1\,d\omega=1.
\]
Thus every such $g_h$ belongs to $S_f$. In particular, $g_0$ and $g_1$ are
distinct elements of $S_f$. By the theorem, $S_f$ is nevertheless weakly
compact.
\end{proof}

\section*{Verification code}

The script \texttt{code/check\_lp\_state\_space\_example.py} checks the
finite-partition version of the last calculation for three choices of $h$.
It verifies that the $L^q$ norm and the pairing with the constant $e_1$
function are both equal to $1$.

\section*{Scope}

This is a partial answer to Problem 3.18. The non-reflexive case remains open
in this packet. The main conclusion is that, for $1<p<\infty$, reflexivity of
$X$ makes weak compactness automatic, and the correct $L^p$ theory cannot be
a direct copy of the source paper's $L^1$ singleton theorem.

\end{document}
